% ================================================================ 3.6 intro: what is DPO
\begin{frame}{A Simpler Path: Direct Preference Optimization (DPO)}
\begin{itemize}\itemsep0.45em\small
    \item \textbf{The pain.} The full pipeline is heavy: Stage 2 trains a
          \emph{separate} reward model, and Stage 3 runs an RL loop juggling several
          models at once (policy, reference, reward model, and, for PPO, a
          critic), which is finicky to tune.
    \item \textbf{The question.} The reward model exists only to score the
          preferences we \emph{already collected}. What if we skip it, and the RL
          loop, and tune the policy \emph{straight from the preference pairs}
          $(y^+, y^-)$?
    \item \textbf{DPO (Rafailov et al., 2023):} exactly that, a single
          supervised-style loss on preference pairs, applied directly to
          \textbf{the policy $\pi_\theta$} (the SFT model we are fine-tuning).
          \empr{No reward model. No RL loop.}
\end{itemize}
\vspace{0.4em}
\begin{center}\small
\setlength{\fboxsep}{5pt}%
\colorbox{boxblue!40}{\parbox{0.92\textwidth}{\centering
\textbf{RLHF:}\; preferences $\to$ train reward model $\to$ RL loop $\to$ policy\\[0.3em]
\textbf{DPO:}\;\; preferences $\to$ one loss $\to$ policy}}
\end{center}
\vspace{0.3em}
{\footnotesize But how can you train toward a reward you never built? The plain-English
idea is on the next slide, then the loss it leads to. The derivation, a page of algebra starting from the RLHF objective, is in Rafailov et al.\ (2023).}
\end{frame}

% ---------------------------------------------------------------- 3.6 intuition (plain English)
\begin{frame}{The DPO Trick, in Plain English}
\textbf{Delete the judge: the chatbot is already its own judge.}
\begin{itemize}\itemsep0.55em
    \item Look at how much \textbf{more likely} the current chatbot is to give an
          answer than the \emph{original} (starting) model was. If it has become much
          more likely to say answer A, that \emph{is} the chatbot ``rating A
          highly.'' \empr{That ratio, ``how much more do I favour this than I used
          to'', is the score.}
    \item So DPO throws away the separate judge \emph{and} the practice loop, and
          trains the chatbot directly so that, for each comparison, \textbf{its own
          favouring of the preferred answer beats its favouring of the rejected
          one.} One model, one training step, no loop.
    \item Rafailov et al.\ make this exact. The next slide shows the loss, and why
          their paper is titled \empr{``your LM is secretly a reward model.''}
\end{itemize}
\end{frame}

